% Lecture 12 - Page 3 (3-Jan-1995).
% Fragment for % Lecture 12 - Page 3 (3-Jan-1995).
% Fragment for % Lecture 12 - Page 3 (3-Jan-1995).
% Fragment for % Lecture 12 - Page 3 (3-Jan-1995).
% Fragment for \input{courses/pages/curs12-pagina3-en}.
% Faithful translation in manuscript order; no preamble or forced page break.

\subsection*{Lecture 12 - Page 3}

\noindent\underline{The free monoid generated by a set}

\bigskip
\noindent Let $A\neq\varnothing$, $1\notin A$ ($A$ is called an \underline{alphabet},
 and the elements of $A$ are called \underline{letters}).

\bigskip
\noindent
\[
A^{+}:=\bigl\{u\mid u:\mathbb{N}^{*}\longrightarrow A\cup\{1\},\ \exists\, n\geq 1
\text{ such that }
\begin{array}{ll}
 u(i)\in A, & (\forall)\ i=1,\ldots,n,\\
 u(i)=1, & (\forall)\ i>n
\end{array}
\bigr\}.
\]

\noindent $u=(x_1,x_2,\ldots,x_m,1,\ldots)$,
\qquad $v\in A^{+}$, $v=(y_1,\ldots,y_n,1,\ldots)$.

\medskip
\noindent We have
\[
u=v\quad\Longleftrightarrow\quad m=n\ \text{and}\ x_i=y_i
\quad(\forall)\ i=1,\ldots,m=n.
\]

\bigskip
\noindent $A^{+}\ni u,v$,
\[
u\cdot v\ \overset{\mathrm{def}}{=}
(x_1,\ldots,x_m,y_1,\ldots,y_n,1,\ldots).
\]
\noindent\hspace*{3em}\textit{concatenation}\hspace{5em}\textit{juxtaposition}
% These two labels are written beneath the operation and its result in the manuscript.

\bigskip
\noindent $(A^{+},\cdot)$ is a semigroup (associative).

\bigskip
\noindent $A^{*}=A^{+}\cup\{1\}$.

\noindent Extend the operation ``$\cdot$'' to $A^{*}$:
\[
u\cdot 1=1u=1,\qquad (\forall)\ u\in A, \qquad 1\cdot 1=1.
\]
% EDITORIAL CHECK: The manuscript appears to end the first equation with "=1".
% For an identity element, the usual equality would be "=u". This is
% deliberately NOT silently amended. "u in A" also appears in the original,
% although the extension would normally be stated for u in A^*.

\bigskip
\noindent\underline{Remark.}\quad $(A^{*},\cdot,1)$ is a monoid,
called the \underline{free monoid generated by $A$}.

\medskip
\noindent $A\subseteq A^{*}$.
% The original depicts this inclusion by a brace beneath A^*.

\noindent
\[
\zeta_A:A\longrightarrow A^{+},\qquad
\zeta_A(x)=(x,1,1,\ldots),\qquad\text{injective map}.
\]
% The map symbol looks like zeta_A; check against the original in case it is iota_A.

\medskip
\noindent\underline{Remark.}\quad
$u=(x_1,x_2,\ldots,x_m,1,1,\ldots)
=\zeta_A(x_1)\zeta_A(x_2)\cdots\zeta_A(x_m)$.

\medskip
\noindent The elements of $A^{+}$ are called \underline{words over $A$};
$m$ is called the length of the word.

\bigskip
\noindent $(A^{*},\cdot,1)$ is a ``strongly'' non-commutative monoid:
\[
x,y\in A,\quad x\neq y\quad\Longrightarrow\quad xy\neq yx.
\]

\bigskip
\noindent\fbox{T.}\quad
\underline{Universal property of the free semigroup generated by a set}

\medskip
\noindent Let $A\neq\varnothing$,
\[A\xrightarrow{\ \zeta_A\ } A^{+}.\]

\noindent For every semigroup $(S,\cdot)$ and every set map
$\zeta:A\longrightarrow S$,
\[
(\exists!)\ f:A^{+}\longrightarrow S\quad\text{a semigroup homomorphism}
\]
\noindent such that $f\circ\zeta_A=\zeta$.

\[
\begin{tikzcd}[column sep=large,row sep=large]
A \arrow[r,"\zeta_A"] \arrow[dr,"\zeta"'] & A^{+} \arrow[d,dashed,"f"] \\
& S
\end{tikzcd}
\]
% The manuscript states the theorem here; no proof is supplied on this page.
.
% Faithful translation in manuscript order; no preamble or forced page break.

\subsection*{Lecture 12 - Page 3}

\noindent\underline{The free monoid generated by a set}

\bigskip
\noindent Let $A\neq\varnothing$, $1\notin A$ ($A$ is called an \underline{alphabet},
 and the elements of $A$ are called \underline{letters}).

\bigskip
\noindent
\[
A^{+}:=\bigl\{u\mid u:\mathbb{N}^{*}\longrightarrow A\cup\{1\},\ \exists\, n\geq 1
\text{ such that }
\begin{array}{ll}
 u(i)\in A, & (\forall)\ i=1,\ldots,n,\\
 u(i)=1, & (\forall)\ i>n
\end{array}
\bigr\}.
\]

\noindent $u=(x_1,x_2,\ldots,x_m,1,\ldots)$,
\qquad $v\in A^{+}$, $v=(y_1,\ldots,y_n,1,\ldots)$.

\medskip
\noindent We have
\[
u=v\quad\Longleftrightarrow\quad m=n\ \text{and}\ x_i=y_i
\quad(\forall)\ i=1,\ldots,m=n.
\]

\bigskip
\noindent $A^{+}\ni u,v$,
\[
u\cdot v\ \overset{\mathrm{def}}{=}
(x_1,\ldots,x_m,y_1,\ldots,y_n,1,\ldots).
\]
\noindent\hspace*{3em}\textit{concatenation}\hspace{5em}\textit{juxtaposition}
% These two labels are written beneath the operation and its result in the manuscript.

\bigskip
\noindent $(A^{+},\cdot)$ is a semigroup (associative).

\bigskip
\noindent $A^{*}=A^{+}\cup\{1\}$.

\noindent Extend the operation ``$\cdot$'' to $A^{*}$:
\[
u\cdot 1=1u=1,\qquad (\forall)\ u\in A, \qquad 1\cdot 1=1.
\]
% EDITORIAL CHECK: The manuscript appears to end the first equation with "=1".
% For an identity element, the usual equality would be "=u". This is
% deliberately NOT silently amended. "u in A" also appears in the original,
% although the extension would normally be stated for u in A^*.

\bigskip
\noindent\underline{Remark.}\quad $(A^{*},\cdot,1)$ is a monoid,
called the \underline{free monoid generated by $A$}.

\medskip
\noindent $A\subseteq A^{*}$.
% The original depicts this inclusion by a brace beneath A^*.

\noindent
\[
\zeta_A:A\longrightarrow A^{+},\qquad
\zeta_A(x)=(x,1,1,\ldots),\qquad\text{injective map}.
\]
% The map symbol looks like zeta_A; check against the original in case it is iota_A.

\medskip
\noindent\underline{Remark.}\quad
$u=(x_1,x_2,\ldots,x_m,1,1,\ldots)
=\zeta_A(x_1)\zeta_A(x_2)\cdots\zeta_A(x_m)$.

\medskip
\noindent The elements of $A^{+}$ are called \underline{words over $A$};
$m$ is called the length of the word.

\bigskip
\noindent $(A^{*},\cdot,1)$ is a ``strongly'' non-commutative monoid:
\[
x,y\in A,\quad x\neq y\quad\Longrightarrow\quad xy\neq yx.
\]

\bigskip
\noindent\fbox{T.}\quad
\underline{Universal property of the free semigroup generated by a set}

\medskip
\noindent Let $A\neq\varnothing$,
\[A\xrightarrow{\ \zeta_A\ } A^{+}.\]

\noindent For every semigroup $(S,\cdot)$ and every set map
$\zeta:A\longrightarrow S$,
\[
(\exists!)\ f:A^{+}\longrightarrow S\quad\text{a semigroup homomorphism}
\]
\noindent such that $f\circ\zeta_A=\zeta$.

\[
\begin{tikzcd}[column sep=large,row sep=large]
A \arrow[r,"\zeta_A"] \arrow[dr,"\zeta"'] & A^{+} \arrow[d,dashed,"f"] \\
& S
\end{tikzcd}
\]
% The manuscript states the theorem here; no proof is supplied on this page.
.
% Faithful translation in manuscript order; no preamble or forced page break.

\subsection*{Lecture 12 - Page 3}

\noindent\underline{The free monoid generated by a set}

\bigskip
\noindent Let $A\neq\varnothing$, $1\notin A$ ($A$ is called an \underline{alphabet},
 and the elements of $A$ are called \underline{letters}).

\bigskip
\noindent
\[
A^{+}:=\bigl\{u\mid u:\mathbb{N}^{*}\longrightarrow A\cup\{1\},\ \exists\, n\geq 1
\text{ such that }
\begin{array}{ll}
 u(i)\in A, & (\forall)\ i=1,\ldots,n,\\
 u(i)=1, & (\forall)\ i>n
\end{array}
\bigr\}.
\]

\noindent $u=(x_1,x_2,\ldots,x_m,1,\ldots)$,
\qquad $v\in A^{+}$, $v=(y_1,\ldots,y_n,1,\ldots)$.

\medskip
\noindent We have
\[
u=v\quad\Longleftrightarrow\quad m=n\ \text{and}\ x_i=y_i
\quad(\forall)\ i=1,\ldots,m=n.
\]

\bigskip
\noindent $A^{+}\ni u,v$,
\[
u\cdot v\ \overset{\mathrm{def}}{=}
(x_1,\ldots,x_m,y_1,\ldots,y_n,1,\ldots).
\]
\noindent\hspace*{3em}\textit{concatenation}\hspace{5em}\textit{juxtaposition}
% These two labels are written beneath the operation and its result in the manuscript.

\bigskip
\noindent $(A^{+},\cdot)$ is a semigroup (associative).

\bigskip
\noindent $A^{*}=A^{+}\cup\{1\}$.

\noindent Extend the operation ``$\cdot$'' to $A^{*}$:
\[
u\cdot 1=1u=1,\qquad (\forall)\ u\in A, \qquad 1\cdot 1=1.
\]
% EDITORIAL CHECK: The manuscript appears to end the first equation with "=1".
% For an identity element, the usual equality would be "=u". This is
% deliberately NOT silently amended. "u in A" also appears in the original,
% although the extension would normally be stated for u in A^*.

\bigskip
\noindent\underline{Remark.}\quad $(A^{*},\cdot,1)$ is a monoid,
called the \underline{free monoid generated by $A$}.

\medskip
\noindent $A\subseteq A^{*}$.
% The original depicts this inclusion by a brace beneath A^*.

\noindent
\[
\zeta_A:A\longrightarrow A^{+},\qquad
\zeta_A(x)=(x,1,1,\ldots),\qquad\text{injective map}.
\]
% The map symbol looks like zeta_A; check against the original in case it is iota_A.

\medskip
\noindent\underline{Remark.}\quad
$u=(x_1,x_2,\ldots,x_m,1,1,\ldots)
=\zeta_A(x_1)\zeta_A(x_2)\cdots\zeta_A(x_m)$.

\medskip
\noindent The elements of $A^{+}$ are called \underline{words over $A$};
$m$ is called the length of the word.

\bigskip
\noindent $(A^{*},\cdot,1)$ is a ``strongly'' non-commutative monoid:
\[
x,y\in A,\quad x\neq y\quad\Longrightarrow\quad xy\neq yx.
\]

\bigskip
\noindent\fbox{T.}\quad
\underline{Universal property of the free semigroup generated by a set}

\medskip
\noindent Let $A\neq\varnothing$,
\[A\xrightarrow{\ \zeta_A\ } A^{+}.\]

\noindent For every semigroup $(S,\cdot)$ and every set map
$\zeta:A\longrightarrow S$,
\[
(\exists!)\ f:A^{+}\longrightarrow S\quad\text{a semigroup homomorphism}
\]
\noindent such that $f\circ\zeta_A=\zeta$.

\[
\begin{tikzcd}[column sep=large,row sep=large]
A \arrow[r,"\zeta_A"] \arrow[dr,"\zeta"'] & A^{+} \arrow[d,dashed,"f"] \\
& S
\end{tikzcd}
\]
% The manuscript states the theorem here; no proof is supplied on this page.
.
% Faithful translation in manuscript order; no preamble or forced page break.

\subsection*{Lecture 12 - Page 3}

\noindent\underline{The free monoid generated by a set}

\bigskip
\noindent Let $A\neq\varnothing$, $1\notin A$ ($A$ is called an \underline{alphabet},
 and the elements of $A$ are called \underline{letters}).

\bigskip
\noindent
\[
A^{+}:=\bigl\{u\mid u:\mathbb{N}^{*}\longrightarrow A\cup\{1\},\ \exists\, n\geq 1
\text{ such that }
\begin{array}{ll}
 u(i)\in A, & (\forall)\ i=1,\ldots,n,\\
 u(i)=1, & (\forall)\ i>n
\end{array}
\bigr\}.
\]

\noindent $u=(x_1,x_2,\ldots,x_m,1,\ldots)$,
\qquad $v\in A^{+}$, $v=(y_1,\ldots,y_n,1,\ldots)$.

\medskip
\noindent We have
\[
u=v\quad\Longleftrightarrow\quad m=n\ \text{and}\ x_i=y_i
\quad(\forall)\ i=1,\ldots,m=n.
\]

\bigskip
\noindent $A^{+}\ni u,v$,
\[
u\cdot v\ \overset{\mathrm{def}}{=}
(x_1,\ldots,x_m,y_1,\ldots,y_n,1,\ldots).
\]
\noindent\hspace*{3em}\textit{concatenation}\hspace{5em}\textit{juxtaposition}
% These two labels are written beneath the operation and its result in the manuscript.

\bigskip
\noindent $(A^{+},\cdot)$ is a semigroup (associative).

\bigskip
\noindent $A^{*}=A^{+}\cup\{1\}$.

\noindent Extend the operation ``$\cdot$'' to $A^{*}$:
\[
u\cdot 1=1u=1,\qquad (\forall)\ u\in A, \qquad 1\cdot 1=1.
\]
% EDITORIAL CHECK: The manuscript appears to end the first equation with "=1".
% For an identity element, the usual equality would be "=u". This is
% deliberately NOT silently amended. "u in A" also appears in the original,
% although the extension would normally be stated for u in A^*.

\bigskip
\noindent\underline{Remark.}\quad $(A^{*},\cdot,1)$ is a monoid,
called the \underline{free monoid generated by $A$}.

\medskip
\noindent $A\subseteq A^{*}$.
% The original depicts this inclusion by a brace beneath A^*.

\noindent
\[
\zeta_A:A\longrightarrow A^{+},\qquad
\zeta_A(x)=(x,1,1,\ldots),\qquad\text{injective map}.
\]
% The map symbol looks like zeta_A; check against the original in case it is iota_A.

\medskip
\noindent\underline{Remark.}\quad
$u=(x_1,x_2,\ldots,x_m,1,1,\ldots)
=\zeta_A(x_1)\zeta_A(x_2)\cdots\zeta_A(x_m)$.

\medskip
\noindent The elements of $A^{+}$ are called \underline{words over $A$};
$m$ is called the length of the word.

\bigskip
\noindent $(A^{*},\cdot,1)$ is a ``strongly'' non-commutative monoid:
\[
x,y\in A,\quad x\neq y\quad\Longrightarrow\quad xy\neq yx.
\]

\bigskip
\noindent\fbox{T.}\quad
\underline{Universal property of the free semigroup generated by a set}

\medskip
\noindent Let $A\neq\varnothing$,
\[A\xrightarrow{\ \zeta_A\ } A^{+}.\]

\noindent For every semigroup $(S,\cdot)$ and every set map
$\zeta:A\longrightarrow S$,
\[
(\exists!)\ f:A^{+}\longrightarrow S\quad\text{a semigroup homomorphism}
\]
\noindent such that $f\circ\zeta_A=\zeta$.

\[
\begin{tikzcd}[column sep=large,row sep=large]
A \arrow[r,"\zeta_A"] \arrow[dr,"\zeta"'] & A^{+} \arrow[d,dashed,"f"] \\
& S
\end{tikzcd}
\]
% The manuscript states the theorem here; no proof is supplied on this page.
